\newpage
\section{Architettura di Sistema}

L'architettura software è stata ingegnerizzata seguendo un paradigma \textbf{stratificato} (\textit{layered architecture}), finalizzato a massimizzare la portabilità del codice e a garantire un netto disaccoppiamento tra l'astrazione dell'hardware e la logica applicativa. Tale approccio permette di gestire la complessità derivante dalla coesistenza di task asincroni e requisiti real-time.

Come illustrato nel diagramma di Fig. \ref{fig:system_architecture}, il sistema è organizzato in tre domini logici principali, ognuno con responsabilità e vincoli temporali distinti.

% --- Inizio Diagramma TikZ (mantenuto come da te fornito) ---
\begin{figure}[H]
\centering
\resizebox{0.9\textwidth}{!}{ 
\begin{tikzpicture}[node distance=2cm, auto,
    % Stile per i blocchi software (Rettangoli)
    block/.style={rectangle, draw=black, thick, fill=white, text width=9em, text centered, minimum height=3.5em},
    % Stile per i componenti hardware (Ellissi)
    hw/.style={ellipse, draw=black, thin, fill=white, text width=7em, text centered, minimum height=2.5em, font=\small\itshape},
    % Stile per le linee di flusso
    line/.style={draw, -latex, thick}]

    % --- LIVELLO 3: APPLICATION (TOP) ---
    \node [block] (app) {\textbf{Logic Manager} \\ \footnotesize (State Machine)};

    % --- LIVELLO 2: MIDDLEWARE / DRIVERS (CENTER) ---
    \node [block, below of=app, yshift=-0.5cm] (storage) {LittleFS \\ \footnotesize (Flash FS)};
    \node [block, left=of storage, xshift=-0.5cm] (gps_p) {minmea \\ \footnotesize (NMEA Parser)};
    \node [block, right=of storage, xshift=0.5cm] (ui_p) {Adafruit GFX \\ \footnotesize (SSD1306 Driver)};

    % --- LIVELLO 1: HARDWARE (BOTTOM) ---
    \node [hw, below of=storage, yshift=-0.5cm] (flash) {W25Q16 \\ \footnotesize (SPI Flash)};
    \node [hw, left=of flash, xshift=-0.5cm] (gps_hw) {SAM-M8Q \\ \footnotesize (GNSS Module)};
    \node [hw, right=of flash, xshift=0.5cm] (oled_hw) {OLED 0.96'' \\ \footnotesize (SSD1306)};

    % --- INPUT ESTERNI ---
    \node [hw, above left=of app, xshift=-1cm] (buttons) {User Buttons \\ \footnotesize (GPIO IRQ)};
    \node [hw, above right=of app, xshift=1cm] (pc) {External PC \\ \footnotesize (Python API)};

    % --- CONNESSIONI DATI (DRITTE E SIMMETRICHE) ---
    
    % Flusso GNSS
    \path [line] (gps_hw) -- node[anchor=east, font=\tiny] {UART} (gps_p);
    \path [line] (gps_p) |- (app);
    
    % Flusso Storage
    \path [line] (app) -- (storage);
    \path [line] (storage) -- node[anchor=east, font=\tiny] {QSPI} (flash);
    
    % Flusso Display
    \path [line] (app) -| (ui_p);
    \path [line] (ui_p) -- node[anchor=west, font=\tiny] {I2C} (oled_hw);
    
    % Input/Output Esterni
    \path [line] (buttons) -- (app);
    \path [line, dashed] (pc) -- (app);

    % --- RAGGRUPPAMENTI ESTETICI (OPZIONALI) ---
    \draw[gray, dotted] (-6.5,-1.2) -- (6.5,-1.2);
    \draw[gray, dotted] (-6.5,-3.8) -- (6.5,-3.8);
    
    \node[anchor=west, gray, font=\tiny\scshape] at (-4.0, -0.8) {Application Layer};
    \node[anchor=west, gray, font=\tiny\scshape] at (-4.0, -3.4) {Middleware Layer};
    \node[anchor=west, gray, font=\tiny\scshape] at (-4.0, -5.5) {Hardware Layer};

\end{tikzpicture}
}
\caption{Architettura di sistema: integrazione simmetrica tra driver hardware e logica di controllo.}
\label{fig:system_architecture}
\end{figure}

% --- SOTTOSEZIONE 1 ---
\subsection{Analisi dei Livelli Architetturali}

La separazione logica introdotta non è solo organizzativa, ma funzionale all'ottimizzazione del \textbf{determinismo} del sistema.

\subsubsection*{Livello Hardware: Physical \& Low-Level I/O}
Rappresenta il livello più basso, dove risiedono i componenti fisici e i driver di periferica. Gestisce il condizionamento dei segnali e il campionamento tramite protocolli standard (UART, I2C, QSPI). L'uso del sottosistema \textbf{PIO} (\textit{Programmable I/O}) dell'RP2040 in questo strato permette di sollevare la CPU dai compiti di timing critico, delegando la generazione delle forme d'onda direttamente all'hardware dedicato.

\subsubsection*{Livello Middleware: Astrazione e Driver}
Agisce come ponte tra l'hardware e la logica. Implementa i protocolli di parsing deterministico (\textit{minmea} per le sentenze NMEA) e la gestione resiliente dei dati tramite il filesystem \textit{LittleFS}. Questo livello astrae la complessità della memoria Flash W25Q16, offrendo interfacce di alto livello per l'archiviazione e il recupero dei log.

\subsubsection*{Livello Application: Logica di Controllo}
Costituisce l'intelligenza centrale del dispositivo. Gestisce una macchina a stati finiti (\textbf{FSM}) che coordina le transizioni del sistema (es. acquisizione, sincronizzazione, errore). Questo strato è agnostico rispetto al protocollo fisico, basandosi esclusivamente sulle strutture dati fornite dai livelli sottostanti.

% --- SOTTOSEZIONE 2 ---
\subsection{Integrazione e Flusso Dati}

Questa granularità architettonica garantisce che le operazioni a latenza variabile, come i cicli di cancellazione della memoria Flash gestiti da \textit{LittleFS}, non interferiscano con la precisione temporale dell'acquisizione GNSS o con la reattività degli interrupt utente.